% =====================================================================
% EXTENDED ANALYSIS --- DRAFT (generated by runners/extra_analysis.py from
% results/benchmark.csv; combined TSB-StreamingAD-U + TSB-StreamingAD-M).
% Compile-safe fragment: no preamble; \input at the end of \section{Results}
% in main.tex. Requires siunitx, booktabs, graphicx, natbib, cleveref, and the
% amssymb \checkmark (all loaded by main.tex). Tables are \input from tables/
% (orientation_stability.tex, nemenyi_pairwise.tex, per_domain.tex); figures are
% cd_diagram, nemenyi_signplot, per_domain_heatmap under figures/. Every numeral
% traces to results/tables/, results/figures/, results/orientation_findings.md,
% and results/extended_findings.md.
% =====================================================================

\subsection{Extended analyses}\label{sec:results-extended}
We close the results with three analyses computed entirely from the completed
per-dataset AUROC matrix (no additional runs): an orientation-stability analysis
that quantifies the inversion of the post-hoc cluster and its dependence on the
normalisation regime; a Nemenyi post-hoc test with a critical-difference diagram
that establishes \emph{which} pairwise separations are statistically significant;
and a per-domain breakdown that localises where each detector family wins and
fails.

\subsubsection{Orientation stability: systematic, regime-dependent inversion}\label{sec:results-orientation}
The overall leaderbord (\cref{tab:tsb_main}) shows the post-hoc cluster sitting
below the \num{0.5} line, but a mean can hide whether the inversion is a
consistent per-dataset phenomenon or an average of scattered wins and losses.
\Cref{tab:orientation_stability} resolves this by reporting, for every detector,
the \emph{inverted fraction} --- the share of datasets on which AUROC${}<0.5$ ---
alongside the sign-flipped mean AUROC ($\overline{1-\text{AUROC}}$) that a single
global sign flip would recover, split by normalisation. The inversion is
systematic. Every post-hoc softmax, logit, activation, and gradient detector is
inverted on a large majority of datasets: GradNorm on \SI{72.1}{\percent}, DICE on
\SI{73.1}{\percent}, ReAct on \SI{70.2}{\percent}, energy on \SI{69.6}{\percent},
ODIN on \SI{69.4}{\percent}, SCALE on \SI{68.1}{\percent}, and even the
least-inverted MSP on \SI{58.4}{\percent}. The feature-manifold leaders are the
mirror image, inverted on under a tenth of datasets (Mahalanobis
\SI{9.5}{\percent}, DFM-PCA \SI{9.9}{\percent}). Because the inversion is
consistent rather than random, a global sign flip would recover most of the lost
separation --- lifting GradNorm from \num{0.294} to \num{0.706}, energy from
\num{0.301} to \num{0.699}, and ODIN from \num{0.303} to \num{0.697} --- direct
evidence that the post-hoc scores carry real ID/OOD signal that simply points the
wrong way. We nonetheless do \emph{not} flip the sign in the reported leaderboard,
for the scientific-integrity reasons developed in \cref{sec:disc-fail}: the
orientation is the method's own paper-defined orientation, and post-hoc
sign-fitting is precisely the malpractice the audit was built to prevent.

The inversion is also regime-dependent, confirming the normalisation dichotomy
(\cref{sec:results-norm}) at the level of individual datasets. Under global
normalisation the post-hoc detectors are inverted far more often than under
per-series normalisation: ReAct \SI{79.9}{\percent} (global) versus
\SI{53.4}{\percent} (per-series), GradNorm \SI{79.6}{\percent} versus
\SI{59.1}{\percent}, energy \SI{79.3}{\percent} versus \SI{52.8}{\percent}, and
DICE \SI{79.3}{\percent} versus \SI{62.2}{\percent}. Global normalisation
preserves the amplitude and level gap of the shifted second source, which is
exactly the far-out magnitude that drives the output head into overconfident
inversion; per-series normalisation rescales that gap away and the inversion
partially --- but never fully --- relaxes. The full per-detector, per-regime
breakdown is recorded in \texttt{results/orientation\_findings.md}.

% auto-generated by runners/extra_analysis.py (orientation stability)
\begin{table}[t]\centering
\caption{Orientation-stability analysis across all completed datasets. For every detector: $N$ datasets scored, mean AUROC, the \emph{inverted fraction} (share of datasets with AUROC${}<0.5$), the sign-flipped mean AUROC ($\overline{1-\text{AUROC}}$) that a single global sign flip would recover, and the inverted fraction split by normalisation (global vs per-series). Detectors are ordered by mean AUROC. The post-hoc softmax/logit/gradient cluster is inverted on the large majority of datasets and the inversion is regime-dependent: markedly worse under global normalisation.}
\label{tab:orientation_stability}
\small\begin{tabular}{lrrrrrr}\toprule
Method & $N$ & Mean AUROC & Inv.\ frac. & Flip AUROC & Inv.\ (global) & Inv.\ (per-ser.) \\\midrule
Mahalanobis & 527 & 0.826 & 0.095 & 0.174 & 0.105 & 0.078 \\
DFM-PCA & 515 & 0.812 & 0.099 & 0.188 & 0.114 & 0.073 \\
DiMMAD & 527 & 0.747 & 0.192 & 0.253 & 0.147 & 0.269 \\
InvAD & 527 & 0.740 & 0.211 & 0.260 & 0.159 & 0.301 \\
CatSight & 515 & 0.735 & 0.217 & 0.265 & 0.176 & 0.288 \\
M2N2 & 527 & 0.733 & 0.226 & 0.267 & 0.186 & 0.295 \\
DEEDEE & 527 & 0.700 & 0.220 & 0.300 & 0.281 & 0.114 \\
CODiT & 527 & 0.684 & 0.190 & 0.316 & 0.183 & 0.202 \\
DiffAD & 527 & 0.531 & 0.493 & 0.469 & 0.476 & 0.523 \\
MSP & 527 & 0.370 & 0.584 & 0.630 & 0.632 & 0.503 \\
SRS & 260 & 0.352 & 0.669 & 0.648 & 0.713 & 0.570 \\
DICE & 527 & 0.313 & 0.731 & 0.687 & 0.793 & 0.622 \\
SCALE & 527 & 0.308 & 0.681 & 0.692 & 0.775 & 0.518 \\
ODIN & 527 & 0.303 & 0.694 & 0.697 & 0.787 & 0.534 \\
ReAct & 527 & 0.303 & 0.702 & 0.697 & 0.799 & 0.534 \\
Energy & 527 & 0.301 & 0.696 & 0.699 & 0.793 & 0.528 \\
GradNorm & 527 & 0.294 & 0.721 & 0.706 & 0.796 & 0.591 \\
\bottomrule\end{tabular}\end{table}


\subsubsection{Nemenyi post-hoc and the critical-difference diagram}\label{sec:results-nemenyi}
The Friedman test (\cref{sec:results-stats}) establishes that the detectors are
not interchangeable, but it does not say \emph{which} pairs differ. We therefore
run a Nemenyi post-hoc test on the complete-case AUROC matrix over the \num{16}
detectors excluding SRS (the full univariate-plus-multivariate pool,
$N = \num{515}$ datasets on which all \num{16} produced a score), first
re-confirming the omnibus Friedman result on this exact matrix ($\chi^2 =
\num{2336.59}$, $p < \num{1e-16}$), in agreement with \cref{sec:results-stats}.
\Cref{fig:cd_diagram} is the resulting critical-difference diagram: detectors are
placed on the mean-rank axis (rank~\num{1} best), and a crossbar joins any group
whose mean ranks differ by less than the critical difference (CD${} =
\num{1.016}$ rank units at $\alpha = 0.05$). The picture is unambiguous.
Mahalanobis (mean rank \num{4.35}) and DFM-PCA (\num{4.79}) occupy the front,
followed by a tight feature-manifold band --- InvAD (\num{6.49}), DiMMAD
(\num{6.54}), CatSight (\num{6.64}), DEEDEE (\num{6.66}), M2N2 (\num{6.71}) --- and
CODiT (\num{7.34}); the entire post-hoc cluster is relegated to mean ranks between
\num{10.38} (MSP) and \num{11.33} (GradNorm). The gap between the weakest
top-tier detector (CODiT, \num{7.34}) and the strongest post-hoc detector (MSP,
\num{10.38}) is roughly three rank units --- three times the critical difference
--- so the separation between the two families is decisively significant.

\Cref{tab:nemenyi_pairwise} makes the head-to-head explicit: each of the five
top-ranked feature-manifold detectors significantly outranks \emph{every} post-hoc
softmax, logit, and gradient detector (all pairwise $p < \num{1e-15}$). There is no
post-hoc detector that any top feature detector fails to beat at $\alpha = 0.05$;
the pairwise significance structure is visualised in \cref{fig:signplot}. The
statistical verdict thus matches the descriptive one: the feature-manifold family
does not merely lead on average, it significantly outranks the post-hoc family
detector-by-detector.

\begin{figure}[t]
\centering
\includegraphics[width=0.95\textwidth]{cd_diagram.png}
\caption{Critical-difference diagram (Nemenyi post-hoc, $\alpha = 0.05$) over the
\num{16} detectors excluding SRS on the complete-case pool of $N = \num{515}$
datasets. Detectors are positioned by mean rank (rank~\num{1}, best, at the top);
a crossbar joins detectors whose mean ranks differ by less than the critical
difference (CD${} = \num{1.016}$). The feature-manifold detectors (Mahalanobis
\num{4.35}, DFM-PCA \num{4.79}, then the InvAD--M2N2 band) are cleanly separated
from the post-hoc cluster (mean ranks \numrange{10.38}{11.33}), which is itself a
statistical tie internally.}
\label{fig:cd_diagram}
\end{figure}

% auto-generated by runners/extra_analysis.py (Nemenyi pairwise)
\begin{table}[t]\centering
\caption{Nemenyi post-hoc pairwise significance: whether each of the five top-ranked feature-manifold detectors significantly outranks every post-hoc softmax/logit/gradient detector (complete-case, $k=16$ detectors, $N=515$ datasets). A checkmark (\checkmark) marks $p<0.05$; all listed contrasts are significant. Mean ranks (1${}={}$best) are given in parentheses.}
\label{tab:nemenyi_pairwise}
\small\begin{tabular}{lccccccc}\toprule
Top detector (rank) & MSP & ODIN & Energy & ReAct & DICE & SCALE & GradNorm \\\midrule
Mahalanobis (4.35) & \checkmark & \checkmark & \checkmark & \checkmark & \checkmark & \checkmark & \checkmark \\
DFM-PCA (4.79) & \checkmark & \checkmark & \checkmark & \checkmark & \checkmark & \checkmark & \checkmark \\
InvAD (6.49) & \checkmark & \checkmark & \checkmark & \checkmark & \checkmark & \checkmark & \checkmark \\
DiMMAD (6.54) & \checkmark & \checkmark & \checkmark & \checkmark & \checkmark & \checkmark & \checkmark \\
CatSight (6.64) & \checkmark & \checkmark & \checkmark & \checkmark & \checkmark & \checkmark & \checkmark \\
\bottomrule\end{tabular}\end{table}


\begin{figure}[t]
\centering
\includegraphics[width=0.7\textwidth]{nemenyi_signplot.png}
\caption{Nemenyi pairwise-significance sign-plot over the \num{16} detectors
(ordered by mean rank). A coloured off-diagonal cell marks a significant pairwise
difference at $\alpha = 0.05$; the dark block separating the feature-manifold
detectors (upper-left) from the post-hoc cluster (lower-right) shows that every
cross-family contrast is significant, while the post-hoc detectors are mutually
indistinguishable.}
\label{fig:signplot}
\end{figure}

\subsubsection{Per-domain behaviour and failure cases}\label{sec:results-perdomain}
Parsing the data-source token from each dataset id (e.g.\ MITDB, WSD, GHL, UCR,
YAHOO, SMD) lets us localise where each family wins and fails.
\Cref{tab:per_domain} reports the mean AUROC by domain and detector family, and
\cref{fig:per_domain_heatmap} gives the finer domain-by-detector heatmap. The
headline is remarkably consistent: the feature-manifold family wins in
\emph{every} domain but one, and the post-hoc family is inverted (below \num{0.5})
in every domain but one --- and it is the \emph{same} one. On the largest domains
the feature family is dominant --- WebService-derived WSD (\num{0.897}, $N =
\num{97}$), SMD (\num{0.937}), Exathlon (\num{0.921}), MSL (\num{0.878}), NAB
(\num{0.838}), and IOPS (\num{0.812}) --- and it is near-perfect on the small,
cleanly separable Genesis (\num{1.000}) domain. The cardiac ECG domains are the
family's hardest: MITDB (\num{0.654}), SVDB (\num{0.612}), and LTDB (\num{0.644})
sit lower, which is also where the post-hoc detectors come closest to chance
(MITDB \num{0.488}, SVDB \num{0.507}), consistent with per-series normalisation
being the registry default for these Medical sources.

The single revealing exception is the PSM (pooled server metrics) domain, and it
is a genuine reversal rather than noise: the feature-manifold family
\emph{fails below chance} there (\num{0.294}) while the post-hoc family
\emph{exceeds \num{0.7}} (\num{0.787}) --- the one domain where the two families
swap places. This is the clearest failure case for the winning family and the
clearest unexpected success for the losing one, and it is domain-localised rather
than scattered, which is why the aggregate ranking is unaffected. Beyond PSM, the
per-dataset failure audit (\texttt{results/extended\_findings.md}) finds
\num{23} datasets on which even the best feature-manifold detector falls below
chance --- predominantly near-stationary STABLE controls from YAHOO and SMAP,
where a specificity detector has little genuine shift to separate, together with
the PSM OOD files --- and \num{106} datasets on which \emph{some} post-hoc
detector exceeds \num{0.7}, most of them the same easy GHL/WSD/YAHOO cases where
any monotone score separates the two sources. These pockets do not overturn the
aggregate result: the feature-manifold advantage is a property of the corpus as a
whole, with PSM the one systematic regime where output-head confidence, not
feature geometry, is the better signal.

\begin{figure}[t]
\centering
\includegraphics[width=0.98\textwidth]{per_domain_heatmap.png}
\caption{Mean AUROC by data-source domain (rows, with dataset counts) and detector
(columns, grouped feature-manifold, then ts-specific, then post-hoc). Green marks
strong separation, red marks inversion. The left (feature-manifold) block is green
across almost every domain; the right (post-hoc) block is red across almost every
domain. The PSM row inverts the pattern --- red on the left, green on the right ---
the single domain where the post-hoc family outperforms the feature family.}
\label{fig:per_domain_heatmap}
\end{figure}

% auto-generated by runners/extra_analysis.py (per-domain x family)
\begin{table}[t]\centering
\caption{Mean AUROC by data-source domain (rows, $N$ datasets each) and detector family (columns): feature-manifold (Mahalanobis, DFM-PCA, DiMMAD, InvAD, CatSight, DEEDEE, M2N2), post-hoc (MSP, ODIN, Energy, ReAct, DICE, SCALE, GradNorm), and ts-specific (CODiT, DiffAD, SRS). The data-source token is parsed from each dataset id. The feature-manifold family wins in every domain; the post-hoc family is inverted (below \num{0.5}) in every domain.}
\label{tab:per_domain}
\small\begin{tabular}{lrrrr}\toprule
Domain & $N$ & feature-manifold & post-hoc & ts-specific \\\midrule
WSD & 97 & 0.897 & 0.165 & 0.562 \\
MITDB & 82 & 0.654 & 0.488 & 0.566 \\
GHL & 76 & 0.735 & 0.258 & 0.534 \\
UCR & 58 & 0.769 & 0.316 & 0.579 \\
YAHOO & 43 & 0.762 & 0.350 & 0.473 \\
SMD & 33 & 0.937 & 0.063 & 0.651 \\
SVDB & 27 & 0.612 & 0.507 & 0.549 \\
SMAP & 23 & 0.770 & 0.259 & 0.663 \\
LTDB & 22 & 0.644 & 0.463 & 0.498 \\
PSM & 18 & 0.294 & 0.787 & 0.339 \\
NAB & 10 & 0.838 & 0.269 & 0.534 \\
MSL & 9 & 0.878 & 0.130 & 0.578 \\
Exathlon & 7 & 0.921 & 0.091 & 0.697 \\
Genesis & 6 & 1.000 & 0.000 & 0.996 \\
IOPS & 6 & 0.812 & 0.273 & 0.467 \\
MGAB & 5 & 0.665 & 0.360 & 0.569 \\
Other & 5 & 0.591 & 0.408 & 0.579 \\
\bottomrule\end{tabular}\end{table}


\subsubsection{The PSM reversal: a boundary condition of manifold-distance
scoring}\label{sec:results-psm-boundary}
The PSM reversal is worth dissecting because it exposes the assumption on which
the entire feature-manifold family rests. All \num{18} PSM OOD files share a
single \num{96}-window test set (\num{48} in-distribution, \num{48} shifted;
their per-detector AUROC is bit-identical), each built by concatenating a PSM
Facility series (Source-1, in-distribution) with an appended SMAP Sensor series
(Source-2, the shift). On this set the direction of every score can be read
directly. For Mahalanobis the in-distribution windows have a \emph{larger} mean
whitened distance to the ID centroid than the shifted windows
(\num{70.6} vs \num{38.9}; medians \num{44.3} vs \num{39.1}), and they are far
more dispersed: the ID-to-OOD distance-variance ratio is \num{14.1}
(\num{11.8} for DFM-PCA). The shifted windows are not outliers of the ID cloud
but a compact pocket nested inside it --- \num{94}\% of them fall within the ID
distances' \numrange{5}{95} percentile band, and their entire range
$[\num{15.6},\num{74.0}]$ lies within the ID range $[\num{15.6},\num{211.0}]$.
Because manifold-distance scoring equates ``far from the ID centroid'' with
``anomalous,'' it ranks these nested, low-variance windows as \emph{more} normal
than genuine ID windows, and the score inverts: Mahalanobis AUROC \num{0.350},
DFM-PCA \num{0.434} (cf.\ \num{0.884} and \num{0.881} on the \num{117} non-PSM
Facility files, where the shift does lie outside the manifold).

The post-hoc detectors succeed on exactly the same windows because they read a
different quantity --- proximity to the classifier's decision boundary rather
than radial distance from the ID centroid. On this set the shift genuinely
lowers output-head confidence: every shifted window has higher softmax
uncertainty than the ID median (MSP medians \num{0.000} vs \num{0.035}, all
\num{48} OOD above the ID median) and higher free energy (Energy medians
\num{-10.7} vs \num{-5.8}, again all \num{48} above), giving Energy AUROC
\num{0.852}, MSP \num{0.805}, ODIN \num{0.836}, and ReAct \num{0.834}
(Figure~\ref{fig:psm_reversal}). The mechanism is thus a clean boundary
condition rather than a stochastic fluke: the feature-manifold approach assumes
the shift lies \emph{outside} the ID manifold, and PSM violates that assumption
with a nested, lower-variance sub-regime, which is precisely the geometry that
inverts distance-based scoring while leaving confidence-based scoring intact. We
state the diagnosis at the level at which it is measured --- the detectors' own
ID-whitened score space, where the inversion is unambiguous; we do not claim it
reduces to a simple raw-signal variance gap, since in raw multivariate space the
appended Source-2 segment is not uniformly lower-variance. What the data show
cleanly is that, in the metric each detector actually uses, PSM's shift is
interior to the ID distances and exterior to the ID confidence region, which is
enough to reverse the two families.

\begin{figure}[t]
\centering
\includegraphics[width=0.85\textwidth]{psm_reversal_hist.png}
\caption{PSM reversal on the shared \num{96}-window OOD test set (representative
file \texttt{OOD\_017\_\ldots\_PSM}). \emph{Left:} Mahalanobis feature distance
--- shifted (OOD) windows are nested \emph{inside} and less dispersed than the
in-distribution (ID) windows, so the distance score inverts (AUROC \num{0.350}).
\emph{Right:} post-hoc Energy --- the same shifted windows sit at higher energy
(lower confidence) than ID windows, so the score ranks correctly (AUROC
\num{0.852}). Higher is presumed-OOD in both panels.}
\label{fig:psm_reversal}
\end{figure}

\subsubsection{Online streaming: incremental vs.\ fit-once
covariance}\label{sec:results-online}
Every result above scores the winning Mahalanobis detector with its covariance
fit \emph{once} on the Source-1 training windows and then frozen. A natural
deployment question is whether the same detector can instead \emph{self-update}
its ID model online as the stream arrives --- never revisiting the training set
--- which would remove the fit-once dependence on a clean warm-up region. We test
this directly on the temporally-ordered TSB-StreamingAD-U evaluation stream
(\texttt{load\_tsb(ordered\_eval=True)}), comparing two variants that share an
identical warm-start: \emph{batch}, the tied covariance fit once and frozen; and
\emph{online}, the same warm-start followed by an unsupervised, score-then-update
mean/covariance refresh after \emph{each} streamed window with exponential
forgetting (decay \num{0.999}). Every numeral below traces to
\texttt{results/online\_incremental.csv} (seed \num{42}, \num{29} datasets across
DRIFT/OOD/STABLE).

The online variant \emph{substantially underperforms} the fit-once batch detector
(\cref{tab:online_incremental}, \cref{fig:online_vs_batch}). Mean per-window AUROC
falls from \num{0.694} (batch) to \num{0.428} (online), a drop of $\Delta =
\num{-0.266}$, and the collapse is uniform across categories: DRIFT
\num{0.681}$\to$\num{0.448}, OOD \num{0.766}$\to$\num{0.479}, and STABLE
\num{0.628}$\to$\num{0.349}. The threshold-aware metrics move the same way ---
mean AUPR \num{0.698}$\to$\num{0.540} and FPR@95 worsens from \num{0.683} to
\num{0.875}. The effect is not an averaging artifact of a few datasets: across the
\num{29} datasets the online variant \emph{wins on \num{1}, ties on \num{2}, and
loses on \num{26}} (per-dataset AUROC, tolerance $\pm\num{0.005}$).

The mechanism is intrinsic to unsupervised self-updating. Because the online
covariance ingests \emph{every} window it scores --- including the anomalous
Source-2 windows, which carry no label to exclude them --- the running ID Gaussian
gradually absorbs the shifted regime into its own support. As the covariance
inflates along the shift directions, the Mahalanobis distance of subsequent
anomalous windows shrinks toward the ID range, and the very separability the
fit-once estimator preserves is eroded window by window. Self-updating thus
defeats the purpose of the detector: the anomalies it is meant to flag are the
same samples that contaminate its model of normality.

Two coverage caveats apply. The reported cells are single-seed (seed \num{42})
and cover \num{10} files per category (\num{29} datasets after the ordered-stream
filter), so this is a representative-coverage probe rather than a full-corpus
sweep; a larger \num{20}-files-per-category run was OOM-killed before completion
and is not reported. The direction and magnitude of the effect are nonetheless
unambiguous (\num{26} of \num{29} losses, $\Delta = \num{-0.266}$), and they carry
a concrete deployment implication: the ``deploy Mahalanobis'' recommendation holds
specifically for the \emph{fit-once batch} estimator, and a naive self-updating
online covariance is \emph{not} a safe substitute for it.

\begin{table}[t]
\centering
\caption{Batch (fit-once) vs.\ online incremental-covariance Mahalanobis under a true streaming protocol: mean per-window AUROC on the temporally-ordered TSB-StreamingAD-U evaluation stream. $\Delta = \text{online} - \text{batch}$.}
\label{tab:online_incremental}
\begin{tabular}{lrrrr}
\toprule
Category & $n$ & Batch AUROC & Online AUROC & $\Delta$ \\
\midrule
DRIFT & 10 & 0.681 & 0.448 & -0.234 \\
OOD & 10 & 0.766 & 0.479 & -0.287 \\
STABLE & 9 & 0.628 & 0.349 & -0.279 \\
\midrule
\textbf{Overall} & 29 & 0.694 & 0.428 & -0.266 \\
\bottomrule
\end{tabular}
\end{table}


\begin{figure}[t]
\centering
\includegraphics[width=0.85\textwidth]{online_vs_batch.png}
\caption{Batch (fit-once) vs.\ online incremental-covariance Mahalanobis on the
temporally-ordered TSB-StreamingAD-U evaluation stream (seed \num{42}, \num{29}
datasets, online decay \num{0.999}). The online self-updating covariance trails
the fit-once batch detector on mean per-window AUROC in every category
(overall $\Delta = \num{-0.266}$; win/tie/loss $= \numlist{1;2;26}$), because
unsupervised per-window updates absorb anomalous windows into the running
in-distribution Gaussian and collapse the Mahalanobis separation.}
\label{fig:online_vs_batch}
\end{figure}
